\documentclass[10pt]{article}
\usepackage[T1]{fontenc}
\usepackage[utf8]{inputenc}
\usepackage[french]{babel}
\usepackage{amsmath,amssymb,mathrsfs,mathpazo}
\usepackage{geometry}
\geometry{a4paper,margin=2cm}
\usepackage{graphicx,array,tabularx,tabularray,colortbl,multirow,multicol,booktabs,enumitem,pifont,ulem,fancybox,diagbox,cancel}
\usepackage{pstricks,pstricks-add,pst-plot,pst-tree,pst-node,pst-text,pst-func,pst-math,pst-bspline,pst-3dplot,pst-eucl,pst-solides3d}
\let\pscalculate\undefined
\usepackage{pst-bezier}
\usepackage{tikz,pgfplots,tkz-tab}
\usetikzlibrary{arrows,arrows.meta,patterns,calc,positioning,trees,decorations.markings,shapes.geometric,shapes.misc,shapes.arrows,intersections,decorations.pathmorphing,decorations.pathreplacing,backgrounds,shadows,babel}
\pgfplotsset{compat=1.18}
\usepgfplotslibrary{fillbetween}
\usepackage[np]{numprint}
\usepackage{listings,makecell,pas-tableur,twoopt}
\let\barre\undefined
\usepackage[pstricks]{tablvar}
\definecolor{darkgreen}{RGB}{0,100,0}
\providecommand{\red}{\color{red}}
\providecommand{\blue}{\color{blue}}
\providecommand{\green}{\color{green}}
\providecommand{\black}{\color{black}}
\providecommand{\white}{\color{white}}
\providecommand{\gray}{\color{gray}}
\providecommand{\cyan}{\color{cyan}}
\providecommand{\magenta}{\color{magenta}}
\providecommand{\yellow}{\color{yellow}}
\providecommand{\violet}{\color{violet}}
\providecommand{\orange}{\color{orange}}
\providecommand{\euro}{\texteuro}
\providecommand{\EUR}{\texteuro}
\providecommand{\ieme}{\textsuperscript{e}}
\providecommand{\degres}{\ensuremath{^\circ}}
\providecommand{\degree}{\ensuremath{^\circ}}
\providecommand{\vv}[1]{\overrightarrow{#1}}
\providecommand{\Coord}[1]{#1}
\providecommand{\decosix}{}
\providecommand{\starredbullet}{\ensuremath{\bullet}}
\pagestyle{empty}
\setlength{\parindent}{0pt}
\usepackage{adjustbox}\DeclareUnicodeCharacter{2236}{:}\DeclareUnicodeCharacter{2190}{\ensuremath{\leftarrow}}
\providecommand{\R}{\mathbb{R}}
\providecommand{\N}{\mathbb{N}}
\providecommand{\Z}{\mathbb{Z}}
\providecommand{\Q}{\mathbb{Q}}
\providecommand{\D}{\mathbb{D}}
\providecommand{\C}{\mathbb{C}}
\providecommand{\e}{{\mathrm{e}}}
\providecommand{\vect}[1]{\overrightarrow{#1}}
\providecommand{\vectt}[1]{\overrightarrow{\mathrm{#1}}}
\providecommand{\vvt}[1]{\overrightarrow{\mathrm{#1}}}
\providecommand{\np}[2][]{{#2~#1}}
\providecommand{\vv}[1]{{\overrightarrow{#1}}}
\renewcommand{\d}{{\mathrm{d}}}
\providecommand{\degres}{{^\circ}}
\providecommand{\degre}{{^\circ}}
\providecommand{\blue}{\color{blue}}
\providecommand{\red}{\color{red}}
\providecommand{\enskip}{\quad}
\providecommand{\llbracket}{\lbrack\!\lbrack}
\providecommand{\rrbracket}{\rbrack\!\rbrack}
\providecommand{\euro}{{\text{€}}}
\providecommand{\no}{{\text{n°}}}
\providecommand{\interval}[2]{{[#1;#2]}}
\providecommand{\intervFO}[2]{{[#1;#2[}}
\providecommand{\dotfill}{\ldots}

\providecommand{\ds}{\displaystyle}
\providecommand{\textding}[1]{#1}
\providecommand{\Oij}{\(\mathrm{O};\vec{\imath},\vec{\jmath}\)}
\providecommand{\Oijk}{\(\mathrm{O};\vec{\imath},\vec{\jmath},\vec{k}\)}
\providecommand{\Ouv}{\(\mathrm{O};\vec{u},\vec{v}\)}
\providecommand{\intervFF}[2]{\ensuremath{[#1;#2]}}
\renewcommand{\euro}{\ifmmode\text{\texteuro}\else\texteuro\fi}
\DeclareUnicodeCharacter{00B0}{\ifmmode^{\circ}\else\textdegree\fi}
\DeclareUnicodeCharacter{20AC}{\ifmmode\text{\texteuro}\else\texteuro\fi}
% Énoncés : Éducation nationale. Transcription et corrigés : APMEP, Denis Vergès et contributeurs. Découpage et classement : CoopMaths, Sébastien Lozano, Éric Elter et contributeurs. Adaptation de la présentation : Maths Explorer.
% CC BY-SA : https://coopmaths.fr/www/about/
% https://coopmaths.fr/alea/static/bac/2022/tex/bac_2022_09_sujet1_ameriquesud_3.tex
\begin{document}
\noindent\textit{APMEP / CoopMaths. CC BY-SA. Présentation adaptée par Maths Explorer.}\par\medskip
\medskip

Soit $g$ la fonction définie sur l'intervalle $]0~;~+\infty[$ par 

\[g(x) = 1+ x^2[1 - 2 \ln (x)].\]

La fonction $g$ est dérivable sur l'intervalle $]0~;~+\infty[$ et on note $g'$ sa fonction dérivée.

On appelle $\mathcal{C}$ la courbe représentative de la fonction $g$ dans un repère orthonormé du plan.

\medskip

\textbf{PARTIE A}

\medskip

\begin{enumerate}
\item Justifier que $g(\text{e})$ est strictement négatif.
\item Justifier que $\displaystyle\lim_{x \to + \infty} g(x) = - \infty$.
\item 
	\begin{enumerate}
		\item Montrer que, pour tout $x$ appartenant à l'intervalle $]0~;~+\infty[$,\: $g'(x) = -4x \ln (x)$.
		\item Étudier le sens de variation de la fonction $g$ sur l'intervalle $]0~;~+\infty[$ 
		\item Montrer que l'équation $g(x) = 0$ admet une unique solution, notée $\alpha$, sur l'intervalle $[1~;~+\infty[$.
		\item Donner un encadrement de $\alpha$ d'amplitude $10^{-2}$
	\end{enumerate}
\item Déduire de ce qui précède le signe de la fonction $g$ sur l'intervalle $[1~;~+\infty[$.
\end{enumerate}

\medskip

\textbf{PARTIE B}

\medskip

\begin{enumerate}
\item On admet que, pour tout $x$ appartenant à l'intervalle $[1~;~\alpha]$, $g''(x) = - 4[\ln (x) + 1]$.

Justifier que la fonction $g$ est concave sur l'intervalle $[1~;~\alpha]$.
\end{enumerate}

\begin{minipage}{0.48\linewidth}
\begin{enumerate}[resume]
\item Sur la figure ci-contre, A et B sont les points de la
courbe $\mathcal{C}$ d'abscisses respectives 1 et $\alpha$.
	\begin{enumerate}
		\item Déterminer l'équation réduite de la droite (AB).
		\item En déduire que pour tout réel $x$ appartenant à
l'intervalle $[1~;~\alpha]$,\: $g(x) \geqslant \dfrac{- 2}{\alpha - 1} x + \dfrac{2\alpha}{\alpha - 1}$.
	\end{enumerate}
\end{enumerate}
\end{minipage} \hfill
\begin{minipage}{0.48\linewidth}
\begin{center}
\psset{unit=1.3cm,arrowsize=2pt 3}
\begin{pspicture}(-0.2,-0.2)(2.5,3)
\psgrid[gridlabels=0pt,subgriddiv=1,gridwidth=0.15pt]
\psaxes[linewidth=1.25pt,labelFontSize=\scriptstyle]{->}(0,0)(0,0)(2.5,3)
\psplot[plotpoints=1000,linewidth=1.25pt,linecolor=red]{1}{1.9}{1 x ln 2 mul sub x dup mul mul 1 add}
\uput[ul](1,2){\red A}\uput[ur](1.9,0){\red B}\uput[d](1.9,0.05){$\alpha$}
\uput[ur](1.5,1.5){\red $\mathcal{C}$}
\end{pspicture}
\end{center}
\end{minipage}

\bigskip

\end{document}
